% Part: incompleteness
% Chapter: arithmetization-syntax
% Section: introduction

\documentclass[../../../include/open-logic-section]{subfiles}

\begin{document}

\olfileid[pt]{inc}{art}{int}
\olsection{Introdução}

Para relacionar computabilidade e lógica, precisamos de uma maneira de
falar dos objetos da lógica (símbolos, termos, !!{formula}s,
!!{derivation}s), das operações sobre eles e de suas propriedades e
relações de uma forma que permita tratamento computacional. Podemos
fazer isso diretamente, considerando funções e relações computáveis
sobre símbolos, sequências de símbolos e outros objetos construídos a
partir deles. Como os objetos da sintaxe lógica são todos finitos e
construídos a partir de conjuntos !!a{enumerable} de símbolos, isso é
possível para alguns modelos de computação. Outros modelos de
computação---como as funções recursivas---restringem-se, porém, aos
números, às suas relações e às suas funções. Além disso, em última
análise também queremos poder tratar da sintaxe dentro de certas
teorias, especificamente em teorias formuladas na linguagem da
aritmética. Nesses casos, é necessário \emph{aritmetizar} a sintaxe,
O.s.\@, representar objetos sintáticos, operações sobre eles e suas
relações, respectivamente, como números, funções aritméticas e relações
aritméticas. A ideia, que remonta a Leibniz, é atribuir números aos
objetos sintáticos.

É relativamente simples atribuir números aos símbolos como seus
``códigos''. Alguns símbolos apresentam uma pequena dificuldade, pois,
p.ex., há infinitas !!{variable}s e até infinitas !!{function}s
de cada aridade~$n$. Mas é claro que se podem atribuir números aos
símbolos sistematicamente, de tal modo que, digamos, $\Obj v_2$ e
$\Obj v_3$ recebam códigos diferentes. Sequências de símbolos (como
termos e !!{formula}s) constituem um desafio maior. Contudo, se
conseguirmos tratar sequências de números de maneira puramente
aritmética (p.ex., pela codificação de sequências mediante
potências de primos), poderemos estender a codificação dos símbolos
individuais à codificação de sequências de símbolos e, depois, a
sequências ou outros arranjos de !!{formula}s, como !!{derivation}s.
Essa codificação ampliada é chamada ``numeração de G\"odel''. A cada
termo, !!{formula} e !!{derivation} atribui-se um número de G\"odel.

Ao codificar sequências de símbolos como sequências de seus códigos e
escolher um sistema de codificação de sequências que possa ser tratado
por funções computáveis, podemos também tratar os números de G\"odel
por funções computáveis. Na prática, todas as funções relevantes serão
recursivas primitivas. Por exemplo, tanto computar o comprimento de uma
sequência quanto computar seu $i$-ésimo elemento a partir do código da
sequência são operações recursivas primitivas. Se o número que codifica
a sequência for, p.ex., o número de G\"odel de !!a{formula}~$!A$,
vemos imediatamente que o comprimento de !!a{formula} e o (código do)
$i$-ésimo símbolo de !!a{formula} também podem ser computados a partir
do número de G\"odel de~$!A$. É um pouco mais difícil provar que, p.ex., a propriedade de ser o número de G\"odel de um termo bem
formado ou de uma !!{derivation} correta é recursiva primitiva. Ainda
assim, isso é possível porque as sequências de interesse (termos,
!!{formula}s, !!{derivation}s) são definidas indutivamente.

Como exemplo, considere a operação de substituição. Se $!A$ é uma
fórmula, $x$ uma variável e $t$ um termo, então $\Subst{!A}{t}{x}$ é o
resultado de substituir toda ocorrência livre de~$x$ em~$!A$ por~$t$.
Suponha agora que tenhamos atribuído números de G\"odel a $!A$, $x$ e
$t$---digamos, respectivamente, $k$, $l$ e $m$. O mesmo esquema atribui
um número de G\"odel a $\Subst{!A}{t}{x}$, digamos,~$n$. Esse
mapeamento---de $k$, $l$ e $m$ para $n$---é o análogo aritmético da
operação de substituição. Quando a operação de substituição leva $!A$,
$x$ e $t$ a $\Subst{!A}{t}{x}$, a função de substituição aritmetizada
leva os números de G\"odel $k$, $l$ e $m$ ao número de G\"odel~$n$.
Veremos que essa função é recursiva primitiva.

A aritmetização da sintaxe não tem apenas interesse abstrato, embora a
percepção de que linguagens como a linguagem da aritmética, que não
dispõem de mecanismos para ``falar sobre'' linguagens, podem afinal
formalizar propriedades complexas de expressões tenha sido
originalmente nada trivial. Daí é apenas um pequeno passo perguntar o
que uma teoria nessa linguagem, como a aritmética de Peano, pode
\emph{provar} acerca de sua própria linguagem (inclusive, p.ex.,
se !!{sentence}s são demonstráveis ou verdadeiras). Isso nos leva aos
famosos teoremas limitativos de G\"odel (sobre a indemonstrabilidade) e
de Tarski (sobre a indefinibilidade da verdade). Mas o recurso de
aritmetizar a sintaxe também é importante para provar alguns resultados
fundamentais da teoria da computabilidade, p.ex., sobre o poder
computacional das teorias ou a relação entre diferentes modelos de
computabilidade. A aritmetização da sintaxe serve de modelo para a
aritmetização de outros objetos e propriedades. Por exemplo, é
igualmente possível aritmetizar configurações e computações (digamos,
de máquinas de Turing). Isso permite simular computações de um modelo
(p.ex., máquinas de Turing) em outro (p.ex., funções
recursivas).

\end{document}
